\begin{proposition}[Determinant factorization]
\label[proposition]{prop:determinant-factorization}
Put
\[
 K_0=\QQ(c,\lambda,e_1,e_2,e_3).
\]
In $K_0[u,v,z]$ one has
\begin{equation}
 \det(M_s)=(u-v)F_8,
\label{eq:determinant-factorization}
\end{equation}
where $F_8$ is homogeneous of degree eight and $u-v\nmid F_8$.
\end{proposition}

\begin{proof}
Each coefficient $q_{ij}$ is homogeneous of degree three in $(u,v,z)$, so
$\det(M_s)$ is homogeneous of degree nine.  The structural factor can be
seen without expanding the determinant.  Put
\[
 h=x-e_1w^2+e_2w-e_3.
\]
The basis above satisfies
\[
 s_0+s_1=h,\qquad s_2=wh.
\]
Hence on the hyperplane $u=v$ one has
\[
 s=h(u+zw).
\]
The section $h\in H^0(C,2\kappa)$ vanishes on the fixed divisor $D$.
Write $\operatorname{div}(h)=D+D'$.  Then
\[
 [\mathcal O_C(D')\otimes\kappa^{-1}]=-x,
 \qquad
 \Nm_\pi\mathcal O_C(D')=\eta^3.
\]
Thus $\pi_*D'$ is a line section of $E$, and the norm of $h$ is, up to a
nonzero scalar, the product of the fixed norm line $Y-x$ and this second
line.  Moreover
\[
 \Nm_\pi(u+zw)=u^3+z^3Y.
\]
After dividing by the fixed line, $Q_s$ is therefore a product of two lines
whenever $u=v$.  Its symmetric matrix has rank at most two, and
$\det(M_s)$ vanishes identically on $u=v$.  Thus $u-v$ divides the
determinant, and the quotient is homogeneous of degree eight.

It remains to show that this factor is simple.  Divisibility of $F_8$ by the
monic linear form $u-v$ is equivalent to the identity
$F_8(v,v,z)=0$ over $K_0$, so one specialization in characteristic zero at
which $F_8(v,v,z)$ is nonzero suffices.  Set
\[
 (c,\lambda,e_1,e_2,e_3)=(1,1,1,0,2),
 \qquad (u,v,z)=(1+t,1,1).
\]
The normalized curve is
\[
 Y^2+2Y=x^3-6x^2+14x-8.
\]
The six conic coefficients and the exact division by $Y-x$ are written in
Appendix~A.1.  Substitution into the determinant gives
\begin{equation}
 4\det M_{1+t,1,1}=2t(24t^3+56t^2+33t+6).
\label{eq:characteristic-zero-simple-factor}
\end{equation}
Hence $4F_8(1,1,1)=12\ne0$ in $\QQ$, so $u-v\nmid F_8$.
\end{proof}

\begin{lemma}[Geometric integrality of the residual octic]
\label[lemma]{lem:residual-octic-integral}
The projective curve
\[
 \mathscr R=V(F_8)\subset\PP^2_{K_0}
\]
is geometrically integral.
\end{lemma}

\begin{proof}
Reduce the normalized specialization
\[
 (c,\lambda,e_1,e_2,e_3)=(1,1,1,0,2)
\]
modulo $13$.  The corresponding short-Weierstrass curve is
\[
 y^2=x^3+2x+5,
\]
whose discriminant is nonzero; the branch cubic obtained from $y=1$ is also
separable.  Multiplying $F_8$ by the nonzero scalar $4$, set
\[
 H(u,v)=4F_8(u,v,1)
\]
at this specialization.  It has total degree eight and $v$-degree eight,
with constant leading coefficient $7$.  At $u=4$ its monic slice is
\[
 \begin{aligned}
 h(v)={}&v^8-3v^7-2v^6+3v^5+4v^4+4v^3\\
       &-3v^2-5v-2\in\mathbf F_{13}[v].
 \end{aligned}
\]
Appendix~A.2 records the Frobenius reductions
\[
 v^{13^8}\equiv v\pmod h,
 \qquad
 \gcd(h,v^{13^4}-v)=1,
\]
together with an explicit B\'ezout identity for the second relation.  The
standard factor-degree criterion then makes $h$ irreducible over
$\mathbf F_{13}$.
If $H$ factored in $\mathbf F_{13}[u,v]$, the constant leading coefficient in
$v$ would force the leading $v$-coefficients of the factors to be units.  A
positive-$v$-degree factor would retain its degree after $u=4$, while a
$v$-degree-zero factor would itself be a unit.  Both possibilities contradict
the irreducibility of $h$.  Hence $H$ is irreducible.

Let $F_{8,\mathrm{sp}}(u,v,z)$ denote the homogeneous specialization of
$F_8$.  Its dehomogenization at $z=1$ is $H/4$, and the nonzero $v^8$ term
shows that $z\nmid F_{8,\mathrm{sp}}$.  If
$F_{8,\mathrm{sp}}=AB$ were a nontrivial homogeneous factorization, then
neither $A(u,v,1)$ nor $B(u,v,1)$ could be a unit: a positive-degree
homogeneous polynomial whose dehomogenization is constant is a scalar
multiple of a power of $z$, which would force $z$ to divide
$F_{8,\mathrm{sp}}$.  Setting $z=1$ would therefore give a nontrivial
factorization of $H$, a contradiction.  Thus the projective special fibre is
irreducible.

The affine curve $H=0$ contains the rational point $(u,v)=(1,6)$, where
\[
 (\partial H/\partial u,\partial H/\partial v)=(2,0).
\]
Thus the point is smooth.  Since $\mathbf F_{13}$ is perfect, the irreducible
special fibre is geometrically reduced.  If it were not geometrically
irreducible, Frobenius would permute its geometric components transitively.
A rational point would then lie on every component and could not be a smooth
point.  The special fibre is therefore geometrically integral.

Spread the coefficients over a localization $R_0=\mathbf Z[1/N]$ with
$13\nmid N$, and choose an affine open
\[
 \mathscr U\subset\mathbb A^5_{R_0}
 =\operatorname{Spec}R_0[c,\lambda,e_1,e_2,e_3]
\]
containing the displayed characteristic-$13$ point on which the required
denominators and one coefficient of $F_8$ are units.  After multiplying by a
unit, $F_8$ defines a flat projective relative Cartier divisor of degree eight
in $\PP^2_{\mathscr U}$.  Geometric integrality is open in such a family
\cite[IV$_3$, Th\'eor\`eme~12.2.4\textup{(viii)}]{EGA4}.  Since the displayed
special fibre is geometrically integral, so is the generic fibre
$\mathscr R$.
\end{proof}

\begin{lemma}[Normalized and full function fields]
\label[lemma]{lem:normalized-function-field}
Work over $e_1\ne0$ and put
\[
 K_{\mathrm n}=\QQ(\delta,\Lambda,E_2,E_3).
\]
On the affine section chart $v\ne0$, put
\[
 U=u/v,\qquad z_0=z/v,\qquad \zeta=e_1z_0.
\]
Let $\mathscr R_{\mathrm n}$ be obtained by setting $e_1=v=1$ and using
coordinates $(\delta,\Lambda,E_2,E_3;U,\zeta)$.  Under the parameter map of
\cref{lem:dominant-chart}, let
\[
 \mathscr R_{T_{\QQ}}
 =\mathscr R\times_{\operatorname{Spec}K_0}
   \operatorname{Spec}\QQ(T_{\QQ}).
\]
Then
\begin{equation}
 k(\mathscr R)=k(\mathscr R_{\mathrm n})(e_1),
 \qquad
 k(\mathscr R_{T_{\QQ}})
 =k(\mathscr R_{\mathrm n})(e_1,q_0),
 \quad q_0^3=\Lambda^{-1}.
\label{eq:residual-function-fields}
\end{equation}
The curve $\mathscr R_{\mathrm n}$ is geometrically integral.  After base
change to $\CC$, the factor-ordering cover on the full chart is the pullback
of the normalized one through these field extensions.
\end{lemma}

